% !TEX root = ../main.tex
\section{Conclusiones grupales}

En la presente práctica de laboratorio se implementó, verificó formalmente y analizó experimentalmente el montículo de Fibonacci y su método de análisis amortizado mediante la función potencial $\Phi(H) = t(H) + 2m(H)$, contrastando su desempeño con el módulo estándar \texttt{heapq} de Python en operaciones básicas, secuencias intensivas de disminución de clave y el algoritmo de caminos mínimos de Dijkstra.

La batería exhaustiva de pruebas automáticas certificó el cumplimiento del $100\%$ de las invariantes estructurales requeridas: orden de montículo, grados exactos por nodo, consistencia bidireccional en listas circulares, ausencia de marcas en las raíces y unicidad de grados tras la consolidación. Se comprobó experimentalmente que el costo amortizado de las operaciones de inserción ($O(1)$), unión ($O(1)$) y disminución de clave ($O(1)$) absorbe rigurosamente el costo de reestructuraciones profundas, tales como cortes en cascada y enlaces múltiples durante la extracción del mínimo ($O(\log n)$), validando la contabilidad teórica donde el potencial liberado financia el trabajo diferido.

En el plano experimental, los resultados evidenciaron la diferencia fundamental entre complejidad asintótica y rendimiento en tiempo de ejecución. En operaciones ordinarias de inserción y extracción completa, la implementación nativa en lenguaje C de \texttt{heapq} superó ampliamente al montículo de Fibonacci debido a la contigüidad de memoria y localidad de caché frente a la sobrecarga de punteros de objetos dinámicos en CPython. No obstante, al someter las estructuras a una tasa dominante de disminuciones ($q = 10n$), el montículo de Fibonacci superó a \texttt{heapq} (568 ms frente a 767 ms para 200\,000 disminuciones), demostrando la ventaja de su cota constante frente a la acumulación desmedida de entradas obsoletas en montículos binarios perezosos (que multiplicaron por once el tamaño físico en memoria). Finalmente, en Dijkstra ambas colas produjeron distancias mínimas idénticas, ratificando que el montículo de Fibonacci es óptimo teóricamente para grafos densos, aunque en sistemas reales interpretados su adopción debe sopesar las elevadas constantes de implementación.
